% ============ 第2章 参考文献(原书 pp.200-202) ============
\section*{参考文献}\addcontentsline{toc}{section}{参考文献}
\begin{thebibliography}{99}
\item[] 注意．本章主要参考 Duren 的专著 \cite{ch2:11} 写成.另外 Hoffman 的书 \cite{ch2:22} 与 Koosis 的专著 \cite{ch2:23} 也被部分参考.
\bibitem[1]{ch2:1} J. Neuwirth and D. J. Newman, Positive $H^{1/2}$ functions are constant, Proc. Amer. Math. Soc. 18(1967), 958.
\bibitem[2]{ch2:2} H. Helson and D. Sarason, Past and future, Math. Scand. 21(1967), 5--16.
\bibitem[3]{ch2:3} H. L. Royden, The boundary values of analytic and harmonic functions, Math. Zeitschr. 78(1962), 1--24.
\bibitem[4]{ch2:4} 竹之内脩, 函数解析, 近代数学講座 13, 朝倉書店, 東京, 1967.
\bibitem[5]{ch2:5} K. deLeeuw and W. Rudin, Extreme points and extremum problems in $H^{1}$, Pacific J. Math. 8(1958), 467--485.
\bibitem[6]{ch2:6} K. deLeeuw, W. Rudin and J. Wermer, The isometries of some function spaces, Proc. Amer. Math. Soc. 11(1960), 694--698.
\bibitem[7]{ch2:7} K. Yabuta, Some uniqueness theorems for $H^{p}(U^{n})$ functions, Tôhoku Math. J. 24(1972), 353--357.
\bibitem[8]{ch2:8} T. Nakazi, Exposed points and extremal problems in $H^{1}$, J. Funct. Anal. 53(1983), 224--230.
\bibitem[9]{ch2:9} D. Sarason, Exposed points in $H^{1}$, I, in Operator Theory: Advances and Applications, 41(1989), 485--496.
\bibitem[10]{ch2:10} J. Inoue, An example of a non-exposed extreme function on the unit ball of $H^{1}$, Proc. Edinburgh Math. Soc. 37(1993), 47--51.
\bibitem[11]{ch2:11} P. L. K. Duren, Theory of $H^{p}$ Spaces, Academic Press, New York, 1970.
\bibitem[12]{ch2:12} T. Nakazi, Extremal problems in $H^{p}$, J. Austral. Math. Soc. 52(1992), 103--110.
\bibitem[13]{ch2:13} E. Hayashi, The solution sets of extremal problems in $H^{1}$, Proc. Amer. Math. Soc. 93(1985), 690--696.
\bibitem[14]{ch2:14} G. Pick, Über die Beschränkungen analytische Funktionen, welche durch vorgegebene Funktionswerte bewirkt werden, Math. Ann. 77(1916), 7--23.
\bibitem[15]{ch2:15} D. Sarason, Generalized interpolation in $H^{\infty}$, Trans. Amer. Math. Soc. 127(1967), 179--203.
\bibitem[16]{ch2:16} L. Carleson, An interpolations problem by bounded analytic functions, Amer. J. Math. 80(1958), 921--930.
\bibitem[17]{ch2:17} H. S. Shapiro and A. Shields, On some interpolation problems for analytic functions, Amer. J. Math. 83(1961), 513--532.
\bibitem[18]{ch2:18} V. Kabaila, Interpolation sequences for the $H_{p}$ classes in the case p < 1, Litovsk. Mat. Sb. 3(1963), 141--147 (in Russian).
\bibitem[19]{ch2:19} W. K. Hayman, Interpolation by bounded functions, Ann. Inst. Fourier (Grenoble) 8(1959), 277--290.
\bibitem[20]{ch2:20} D. J. Newman, Interpolation in $H^{\infty}$, Trans. Amer. Math. Soc. 92(1959), 501--507.
\bibitem[21]{ch2:21} J. Garnett, Harmonic interpolating sequences, $L^{p}$ and BMO, Ann. Inst. Fourier (Grenoble) 28, 215--228.
\bibitem[22]{ch2:22} K. Hoffman, Banach Spaces of Analytic Functions, Prentice-Hall, Englewood Cliffs, New Jersey, 1962.
\bibitem[23]{ch2:23} P. Koosis, Introduction to $H_{p}$ Spaces, London Math. Society Lecture Notes Series, 40. Cambridge Univ. Press, London and New York, 1980.
\end{thebibliography}